% !TeX root = ../papers/nc-p3.tex
\section{非被約な印から反対位置の射影同一視を延ばす}
\label{sec:finite-marked-projection-symmetry}
四辺が有限な標準有限Frobenius候補で全ての中支持が4とする。
共通の長さ20の点スキームを $X$、その各座標を $x_i$ とし、一つの閉点の座標を $p_i$ とする。
三行三列の行列を、位置1と3で $p_1,p_3$ を消す双線型形式として読む。
\begin{theorem}\label{thm:finite-marked-projection-symmetry}
二つの四次元空間 $w(p_0,-,V_2^*,-)$, $w(V_0^*,-,p_2,-)$ が一致することを仮定し、
共通空間を $E$ とする。反対階数8ならこの一致はどの印でも自動である。
次の二条件を課す。
\begin{enumerate}
\item 印以外の各閉点で、$x_1,x_3$ の $p_1,p_3$ からの射影を $l,r$ とする。
$E$ の中で $l^tZ=Zr=0$ を満たす空間は一次元である。
\item 印の接空間次元は高々2で、各非零接方向を二つの射影座標の微分で
$l,r$ に写したときにも同じ一次元条件が成り立つ。
ここで微分は $p_i$ のアフィン近傍でその点を原点とした三座標について取る。
\end{enumerate}
このとき唯一の線型同定 $T:V_0^*\to V_2^*$,
$w(x,-,p_2,-)=w(p_0,-,Tx,-)$ は、$X$ 上で反対位置への二移動を射影的に与える。
窓はこの位置について対称になる。他方でも同じ条件を満たす印があれば両対称となり、
定理\ref{thm:bisymmetric-two-length-test}によって全次数の必要十分判定が得られる。
二つの印は異なってよい。

特に $E$ の行列式三次曲面が滑らかなら二つの一次元条件が成り立つ。
より一般に、共通の左右核を持つ行列の二次元鉛筆が $E$ 内にないことでも十分である。
\end{theorem}
\begin{proof}
二つの部分切断はどちらも四文字空間から $E$ への同型である。
独立な二つの側の係数を反対切断の行と列に分けて考える。
偶位置の一つが印の行は、奇位置の一つが印の列を全て消し、その逆も成り立つ。
行側の像と列側の像は各々次元4以上で、その和は反対階数以下である
（独立行を選んだ後、それらを消す列を選ぶ）。階数8なら両側で次元4なので一致を得る。

局所的に各点座標の単元倍を選ぶ。$A=w(p_0,-,x_2,-)$ と $B=w(x_0,-,p_2,-)$ は
いずれも $E$ に値を取り、閉点では零にならない。
両方とも、左右から同じ $x_1,x_3$ で縮約すると零である。これは四つの隣接点対の関係そのもので、
$X$ 上の恒等式である。$E$ の全元は $p_1,p_3$ をも消すため、これらの縮約には三座標 $l,r$ を用いられる。
印以外では六本の線型条件の階数は3であり、$A$ がその核の直和因子を張る。
従って $B$ は $A$ の単元倍である。これは局所環が非被約でも成り立つ。

印の局所環を $S$、極大イデアルを $\mathfrak m$ とする。$A,B$ の剰余は単元倍を除いて
$M=w(p_0,-,p_2,-)\ne0$ である。$\lambda(M)=1$ となる $E$ の汎関数を取り、
$Z=B-\lambda(B)A/\lambda(A)$ と置く。$Z$ は三次元の $N=\ker\lambda$ に値を持ち、
三座標 $l,r\in\mathfrak m^{\oplus3}$ による六つの縮約を全て消す。
接方向上の一次の $l,r$ は $M$ を消す（$l^tA=Ar=0$ を一次で読む）。
仮定から $N$ 上の $6\times3$ 縮約行列の一次項は全非零接方向で階数3である。

この条件が $\mathfrak m^3Z=0$ を与えることを示す。
接次元2なら、一次項の転置により接方向の $\mathbb P^1$ 上に全射
$\mathcal O^6\to\mathcal O(1)^3$ がある。核は三つの
$\mathcal O(-a_i)$ の直和で、$a_i\ge0,\ \sum a_i=3$。
従って2だけひねれば核の一次コホモロジーが消え、行の二次式倍は
三個の任意の二変数三次式を生成する。接次元1では非零方向での定数係数の階数3から同じ主張を得る。
接次元0なら自明である。局所環の次数環は対称代数の商なので、
実際の行の $\mathfrak m^2$ 倍は $\mathfrak m^3 S^{\oplus3}$ を $\mathfrak m^4$ を法として張る。
さらに $\mathfrak m$ の冪を掛けて冪零性を使えば $\mathfrak m^3 S^{\oplus3}$ の全体が行加群に入る。
これらの行は全て $Z$ を消すので主張が従う。

最後にこの三乗の誤差を取り除く。$D=H^0(\mathcal O_X)$ とし、
$J$ を印で $\mathfrak m^3$、他の因子で単位イデアルとする。
隣接二位置の座標積 $C=U_0U_1\subset D$ は定理\ref{thm:finite-square-chart}の完全跡対合で
自分自身と直交する。実際 $C^2=U_0^2U_1^2$ をこの跡が消す。
さらに $C\to D/J$ は全射である。各位置で単元座標を1に自明化すると、両位置の
一次座標差は閉埋込みのためそれぞれ $\mathfrak m/\mathfrak m^2$ を張り、
両者の積は $\mathfrak m^2/\mathfrak m^3$ を張る。従って
$C\cap\operatorname{ann}(J)=0$ である。実際この共通部分の元は
$J$ と $C$ の両方に跡直交し、$C+J=D$ である。
第二辺の任意の関係を $(x_1,T x_0)$ で評価した値は $C$ に入り、
印以外で零、印では $\mathfrak m^3$ に消される。後者は $T x_0$ と $x_2$ が
単元倍を除いて $Z$ の線型成分だけ異なることによる。
よってこの評価値は全て零である。第二辺の一意なグラフにより $T$ は $X$ 上で二移動を与える。
六本の関係は有限完全交叉の $(1,1)$ 消去式を全て張るので最初の二辺は反転となり、
命題\ref{prop:finite-flip-obstructions}の交代三次交叉の消失からテンソルの対称性が従う。

二つの縮約を消す元は各閉点では $A$ により、接方向では $M$ により既に一つある。
二元あれば、それらが張る射影直線では基底を取って全行列の最終行と最終列が零になる。
行列式が滑らかな曲面なら階数1以下の点は存在しないが、この直線の二行二列の行列式は
二変数二次式であり零点を持つ。従って滑らかな場合にもこのような直線はない。
\end{proof}
